\documentclass[aspectratio=169,11pt]{beamer}
\usepackage{fathomnet-course}
% generated by write_numbers.py from results/reference_results.json - do not edit
\newcommand{\RefCrops}{23{,}664}
\newcommand{\RefDives}{1{,}465}
\newcommand{\RefTopTenDives}{21\,\%}
\newcommand{\RefTwoTrain}{444}
\newcommand{\RefTwoTest}{148}
\newcommand{\RefThr}{0.05}
\newcommand{\RefThrTrain}{98.4\,\%}
\newcommand{\RefThrTest}{98.0\,\%}
\newcommand{\RefCSPhylumMajority}{30.4\,\%}
\newcommand{\RefCSPhylumCentroid}{26.4\,\%}
\newcommand{\RefCSPhylumKnn}{55.0\,\%}
\newcommand{\RefCSPhylumTree}{50.3\,\%}
\newcommand{\RefCSLabelMajority}{1.3\,\%}
\newcommand{\RefCSLabelCentroid}{15.9\,\%}
\newcommand{\RefCSLabelKnn}{33.6\,\%}
\newcommand{\RefCSLabelTree}{23.7\,\%}
\newcommand{\RefCSLabelBest}{33.6\,\%}
\newcommand{\RefZooBestName}{neural network (MLP)}
\newcommand{\RefZooBest}{64.6\,\%}
\newcommand{\RefZooBestPixels}{46.3\,\%}
\newcommand{\RefZooBestColour}{48.4\,\%}
\newcommand{\RefZooRows}{majority baseline & 1.3\,\% & 1.3\,\% & $<$1\,s \\
nearest centroid & 39.9\,\% & 37.9\,\% & 2\,s \\
5-nearest neighbours & 69.9\,\% & 53.4\,\% & $<$1\,s \\
decision tree & 100.0\,\% & 20.5\,\% & 48\,s \\
random forest & 100.0\,\% & 56.1\,\% & 93\,s \\
logistic regression & 90.7\,\% & 63.7\,\% & 21\,s \\
neural network (MLP) & 92.4\,\% & 64.6\,\% & 13\,s \\}
\newcommand{\RefRandomAcc}{63.5\,\%}
\newcommand{\RefRandomAccKnn}{53.5\,\%}
\newcommand{\RefTrainLogregRandom}{90.7\,\%}
\newcommand{\RefTrainKnnRandom}{70.0\,\%}
\newcommand{\RefImageAcc}{57.5\,\%}
\newcommand{\RefImageAccKnn}{45.5\,\%}
\newcommand{\RefTrainLogregImage}{91.1\,\%}
\newcommand{\RefTrainKnnImage}{71.0\,\%}
\newcommand{\RefDiveAcc}{49.6\,\%}
\newcommand{\RefDiveAccKnn}{35.5\,\%}
\newcommand{\RefTrainLogregDive}{91.0\,\%}
\newcommand{\RefTrainKnnDive}{71.6\,\%}
\newcommand{\RefSeenDiveRandom}{97\,\%}
\newcommand{\RefRecallRows}{species & 33 & 67\,\% \\
genus & 21 & 69\,\% \\
family & 8 & 64\,\% \\
order & 6 & 56\,\% \\
class & 6 & 44\,\% \\
phylum & 1 & 27\,\% \\
other rank & 4 & 58\,\% \\}
\newcommand{\RefSevKingdom}{55\,\%}
\newcommand{\RefSevNear}{21\,\%}
\newcommand{\RefRandPhylumAcc}{80.0\,\%}
\newcommand{\RefDivePhylumAcc}{72.6\,\%}
\newcommand{\RefDiveFone}{0.475}
\newcommand{\RefDiveAccSingle}{49.8\,\%}

\title{Mini-challenge, recap and exam preparation}
\subtitle{\coursetitle\\[2pt]\small Session 5}
\date{}
\author{}
\begin{document}
\maketitle

\section{The mini-challenge}

\begin{frame}{The task}
\begin{block}{Build the best classifier for a new expedition}
79 labels $\cdot$ test set = dives the classifier has never seen (\texttt{data.split(df, "dive", seed=0)}) $\cdot$ score = \textbf{macro-F1}
\end{block}
\begin{columns}[T]
\column{0.5\textwidth}
\textbf{You may change}
\begin{itemize}\small
  \item features (other CNNs: ResNet-50, EfficientNet; combinations)
  \item classifier and its settings
  \item which training crops you use
  \item Level 5: fine-tune a whole CNN on a GPU
\end{itemize}
\column{0.46\textwidth}
\textbf{The rules}
\begin{enumerate}\small
  \item The test set is fixed.
  \item Choose using a \textbf{validation} split inside the training dives.
  \item Evaluate on the test set \textbf{once}; submit that number.
\end{enumerate}
\end{columns}
\vspace{4pt}
\muted{\small Reference to beat: CNN (ResNet-18) + logistic regression, macro-F1 \RefDiveFone\ (accuracy \RefDiveAccSingle).}
\end{frame}

\begin{frame}{Why a validation set?}
\centering
\begin{tikzpicture}[>=Stealth, font=\small]
  \fill[sea!25] (0,0) rectangle (5.5,0.9);
  \fill[kelp!35] (5.5,0) rectangle (8.6,0.9);
  \fill[coral!35] (8.6,0) rectangle (11.2,0.9);
  \node at (2.75,0.45) {\textbf{train} (dives A)};
  \node at (7.05,0.45) {\textbf{validation} (dives B)};
  \node at (9.9,0.45) {\textbf{test} (dives C)};
  \draw[decorate, decoration={brace, amplitude=5pt}] (0,1.05) -- (8.6,1.05) node[midway, above=5pt] {try 20 ideas here, keep the best};
  \draw[->, thick] (9.9,-0.1) -- (9.9,-0.7) node[below] {use once, at the end};
\end{tikzpicture}

\vspace{12pt}
\begin{alertblock}{If you choose on the test set \dots}
\dots\ you pick the idea that was \emph{lucky} on those particular dives. Your reported score is then too optimistic --
the same trap as testing on the training set, only smaller.
\end{alertblock}
\end{frame}

\begin{frame}{Your poster (A3, or 4 slides)}
\small
\begin{enumerate}
  \item \textbf{Question \& data} -- one sentence each; acknowledge FathomNet / MBARI.
  \item \textbf{Method} -- features and classifier, in words a biologist understands. No code.
  \item \textbf{Results} -- test macro-F1 and accuracy; random vs.\ dive split; one confusion matrix.
  \item \textbf{Mistakes} -- three mistakes with crops and a biological explanation.
  \item \textbf{Verdict} -- would you use it for a survey on a new expedition? What must be checked by hand?
  \item \textbf{What you tried} -- including the ideas that did \emph{not} work.
\end{enumerate}
\vspace{4pt}
\begin{exampleblock}{Graded on}
correct use of train / validation / test $\cdot$ clear figures $\cdot$ biological interpretation $\cdot$ honest verdict. The score itself is only a small bonus.
\end{exampleblock}
\end{frame}

\section{Recap}

\begin{frame}{Recap 1: what a classifier is}
\begin{columns}[T]
\column{0.48\textwidth}
\begin{itemize}\small
  \item \textbf{classifier}: rule $f(\text{features}) = \text{label}$, learned from a \textbf{training set}
  \item \textbf{features}: numbers describing a specimen (characters)
  \item \textbf{decision boundary}: where the predicted label changes
  \item \textbf{baseline}: always the most common label
\end{itemize}
\column{0.48\textwidth}
\small
\begin{tabular}{@{}ll@{}}
\toprule
nearest centroid & closest average \\
$k$-NN & vote of $k$ most similar \\
decision tree & identification key \\
random forest & many keys, vote \\
logistic regression & weighted checklist \\
neural network & checklists of checklists \\
\bottomrule
\end{tabular}
\end{columns}
\end{frame}

\begin{frame}{Recap 2: features, overfitting, testing}
\begin{itemize}
  \item \textbf{Features matter most.} Pixels $<$ colour characters $<$ CNN features (\RefZooBestPixels\ $\rightarrow$ \RefZooBestColour\ $\rightarrow$ \RefZooBest).
  \item \textbf{Transfer learning}: a CNN trained on everyday photos provides good characters for deep-sea animals.
  \item \textbf{Overfitting}: training accuracy rises with flexibility, test accuracy does not. Training accuracy measures memory.
  \item \textbf{Train / validation / test}: learn / choose / report once.
  \item \textbf{Data leakage}: near-identical crops (same dive) on both sides inflate the result: \RefRandomAcc\ $\rightarrow$ \RefDiveAcc.
  \item \textbf{Biological severity}: ask how far apart truth and prediction are in the tree of life.
\end{itemize}
\end{frame}

\section{The final exam}

\begin{frame}{Final exam -- format}
\begin{columns}[T]
\column{0.52\textwidth}
\begin{itemize}\small
  \item written, \textbf{90 minutes}, 100 points
  \item no code; pocket calculator allowed
  \item one handwritten A4 sheet (both sides) allowed
\end{itemize}
\vspace{4pt}
\small
\begin{tabular}{@{}lr@{}}
\toprule
A. Concepts (short answers) & 20 \\
B. Accuracy, precision, recall & 20 \\
C. $k$-NN and a key by hand & 20 \\
D. Interpreting an experiment & 20 \\
E. Mistakes in the tree of life & 10 \\
F. Design a survey classifier & 10 \\
\bottomrule
\end{tabular}
\column{0.44\textwidth}
\begin{block}{Example question}
\small A classifier reaches 99\,\% accuracy on its training crops and 41\,\% on crops from new dives.
Name the phenomenon, give two possible reasons, and say which number belongs in a report.
\end{block}
\end{columns}
\end{frame}

\begin{frame}{How to prepare}
\begin{enumerate}
  \item Re-read the \textbf{\color{coral}Q} answers of your lab notebooks.
  \item Be able to compute by hand: accuracy, recall, precision and the majority baseline from a small confusion matrix; Euclidean distance and a $k$-NN vote; a path through a decision tree.
  \item Be able to explain in two sentences: overfitting, transfer learning, why a dive split, validation vs.\ test.
  \item Be able to judge a mistake with the tree of life.
\end{enumerate}
\end{frame}

\begin{frame}{Outlook: what the professionals add}
\small
\begin{itemize}
  \item \textbf{Fine-tuning} the whole CNN on deep-sea images (Lab 4, Level 5), with \emph{data augmentation} (flips, colour changes).
  \item \textbf{Ensembles}: average several networks (the 3rd-place team used 8 networks with 100 heads each).
  \item \textbf{Taxonomy-aware decisions}: choose the label with the smallest \emph{expected} tree distance instead of the most probable one (5th place).
  \item \textbf{Context}: look at the surroundings, not only the crop (1st and 2nd place) -- like an expert who knows the habitat.
\end{itemize}
\vspace{4pt}
\begin{exampleblock}{Follow-up project (optional)}
``Do a CNN's mistakes follow the tree of life?'' -- run the published winning models on new FathomNet images and analyse their errors by taxonomic rank.
\end{exampleblock}
\end{frame}

\begin{frame}[standout]
Good luck with the challenge!\\[8pt]
\normalsize Next session: poster session \& final exam
\end{frame}
\end{document}
